\documentclass[11pt]{article}

\usepackage[margin=1in]{geometry}
\usepackage{amsmath}
\usepackage{amssymb}
\usepackage{booktabs}
\usepackage{hyperref}

\title{Phase 2A Dynamic FJSP Minimal Model}
\author{FJSP-AGV-RL Research}
\date{June 2026}

\begin{document}
\maketitle

\section{Scope}

Phase 2A extends the verified Phase 1 static flexible job shop scheduling
platform with a minimal dynamic layer. The current scope is deliberately small:
dynamic events, a dynamic state, frozen records, remaining-operation views, and
deterministic rescheduling. AGV routing, transportation time, reinforcement
learning, graph neural networks, operation preemption, and multi-objective
optimization remain outside this phase.

\section{Static Base}

Let $J$ be the set of jobs and let job $j \in J$ contain ordered operations
$O_j = \{0,\ldots,n_j-1\}$. Each operation $(j,o)$ has an eligible machine set
$M_{jo}$ and processing time $p_{jom}$ for each $m \in M_{jo}$. A schedule record
is
\[
  r = (j,o,m,s,e,p),
\]
where $s$ is start time, $e=s+p$ is end time, and $p=p_{jom}$. Phase 2A keeps the
same \texttt{ScheduleRecord} and \texttt{ScheduleResult} representation used by
Phase 1.

\section{Dynamic State}

At rescheduling time $t$, the dynamic state is
\[
  S_t = (I_t, t, R_t, a_t, u_t),
\]
where:

\begin{itemize}
  \item $I_t$ is the current FJSP instance.
  \item $R_t \subseteq J_t$ is the set of released jobs.
  \item $a_t(j)$ is the release time of job $j$.
  \item $u_t(m)$ is the earliest time machine $m$ can start a new future
        operation.
\end{itemize}

The implementation stores $a_t$ as \texttt{job\_release\_times} and $u_t$ as
\texttt{machine\_unavailable\_until}.

\section{Events}

The Phase 2A event set is shown in Table~\ref{tab:events}. Events are applied in
non-decreasing time order by \texttt{apply\_events()}.

\begin{table}[h]
\centering
\begin{tabular}{ll}
\toprule
Event & State effect \\
\midrule
\texttt{MachineBreakdown} & Set $u_t(m)$ to at least the repair time. \\
\texttt{MachineRepair} & Set $u_t(m)$ to the repair event time. \\
\texttt{JobArrival} & Append a new job and set its release time. \\
\texttt{ProcessingTimeChange} & Update one future $p_{jom}$. \\
\bottomrule
\end{tabular}
\caption{Phase 2A dynamic event types.}
\label{tab:events}
\end{table}

\section{Frozen Records}

Given an existing schedule $x$ and current time $t$, Phase 2A defines the frozen
record set as
\[
  F_t =
  \{r \in x : e_r \le t\}
  \cup
  \{r \in x : s_r < t < e_r\}.
\]
The second set freezes in-process operations because Phase 2A does not model
preemption. The helper \texttt{frozen\_records\_at\_time()} implements this rule.

Frozen records must be copied exactly into the rescheduled result:
\[
  (m'_r,s'_r,e'_r,p'_r) = (m_r,s_r,e_r,p_r), \quad r \in F_t.
\]

\section{Remaining Operations}

For each released job $j$, frozen operations must form a route prefix:
\[
  F_{tj} = \{0,\ldots,k_j-1\}.
\]
The remaining operations are then
\[
  Q_t = \{(j,o): j \in R_t,\ o \in \{k_j,\ldots,n_j-1\}\}.
\]
The helper \texttt{remaining\_operations\_from\_frozen()} exposes $Q_t$ for audit
and later repair policies.

\section{Rescheduling Baseline}

The current baseline dispatches all operations in $Q_t$ while respecting:

\begin{align}
  s_{jo} &\ge a_t(j), && \text{job release time,}\\
  s_{jo} &\ge u_t(m), && \text{machine availability for future starts,}\\
  s_{j,o+1} &\ge e_{jo}, && \text{job precedence,}\\
  [s_{jo}, e_{jo}) &\cap [s_{j'o'}, e_{j'o'}) = \emptyset,
    && \text{same-machine non-overlap.}
\end{align}

The implementation currently supports deterministic \texttt{fifo},
\texttt{spt}, and \texttt{eft} rules. The returned result is still a full
\texttt{ScheduleResult}; it is checked by both the Phase 1 static feasibility
checker and the Phase 2A dynamic feasibility checker.

\section{Smoke Verification}

Experiment \texttt{exp\_014\_dynamic\_reschedule\_smoke.py} currently validates
three small scenarios:

\begin{table}[h]
\centering
\begin{tabular}{lrr}
\toprule
Scenario & Makespan & Feasible \\
\midrule
\texttt{dynamic\_breakdown\_1x1} & 12 & PASS \\
\texttt{dynamic\_job\_arrival\_2x1} & 10 & PASS \\
\texttt{dynamic\_multi\_event\_2x2} & 8 & PASS \\
\bottomrule
\end{tabular}
\caption{Current Phase 2A dynamic smoke results.}
\end{table}

\section{Limitations}

The current model is intentionally conservative. It does not interrupt running
operations, add transport resources, optimize AGV paths, simulate stochastic
event streams, or train policies. These restrictions keep Phase 2A small enough
to audit while preserving the Phase 1 legality guarantees.

\end{document}
